\documentclass[11pt,letterpaper]{article}
\usepackage{amsmath}
\usepackage{geometry}
\usepackage{booktabs}
\usepackage{longtable}
\usepackage{hyperref}

\geometry{margin=1.0in}
\hypersetup{colorlinks=true,linkcolor=blue!70!black,urlcolor=blue}

\title{\textbf{RS05 Motor Notes for Bistable Hook Excitation}}
\author{Rolly Robot}
\date{2026}

\begin{document}
\maketitle

\section{Motor Parameters}

\begin{table}[h]
\centering
\caption{RS05 parameters used by the bistable simulations.}
\renewcommand{\arraystretch}{1.25}
\begin{tabular}{p{3.6cm} p{3.2cm} p{6.4cm}}
\toprule
\textbf{Parameter} & \textbf{Value} & \textbf{Notes} \\
\midrule
Rated torque & 1.6\,N$\cdot$m & Continuous rating with heat sinking. \\
Peak torque & 5.5\,N$\cdot$m & Short-duration command limit. \\
No-load output speed & 480\,RPM $\pm10\%$ & Equivalent to $\omega_{\max}\approx50.3$\,rad/s at the output shaft. \\
Torque constant $K_t$ & 0.94\,N$\cdot$m/A$_{\mathrm{rms}}$ & Used to convert Iq current to torque. \\
Gear ratio & 7.75:1 & Planetary reduction. \\
Drive mode & FOC & Field-oriented control. \\
Maximum phase current & 11\,A$_{\mathrm{pk}}$ & The torque command is still limited to the motor's peak torque range. \\
CAN bus rate & 1\,Mbps & CAN 2.0B extended frame. \\
Encoder & 14-bit absolute & Used for snap detection and state feedback. \\
\bottomrule
\end{tabular}
\end{table}

\section{Velocity Limit}

The 480\,RPM value is a shaft angular velocity limit, not a direct frequency
limit. For an oscillating shaft command at frequency $f$ and angular amplitude
$\Theta$,
\begin{equation}
  v_{\mathrm{peak}}=\Theta\,2\pi f \leq \omega_{\max}.
\end{equation}
Therefore
\begin{equation}
  \Theta_{\max}(f)=\frac{\omega_{\max}}{2\pi f}.
\end{equation}
With a lever arm $r$, the maximum linear stroke is
\begin{equation}
  \delta_{\max}(f)=\frac{\omega_{\max}r}{2\pi f}.
\end{equation}

\section{CAN Control Modes}

The RS05 run mode is selected by writing \texttt{run\_mode} at index
\texttt{0x7005} with Communication Type~18, then enabling the motor with
Type~3.

\begin{table}[h]
\centering
\caption{RS05 control modes relevant to sinusoidal excitation.}
\renewcommand{\arraystretch}{1.25}
\begin{tabular}{c p{3.0cm} p{8.5cm}}
\toprule
\textbf{Value} & \textbf{Mode} & \textbf{Use} \\
\midrule
0 & Operation & MIT-style torque/position/velocity command. Set $K_p=K_d=0$ for pure torque feed-forward through $t_{ff}$. \\
3 & Current & Direct Iq command through \texttt{iq\_ref}; no position or speed loop. \\
5 & Position CSP & Cyclic synchronous position; useful only if the position trajectory and velocity limit are acceptable. \\
\bottomrule
\end{tabular}
\end{table}

\section{Streaming A Sine Command}

The motor firmware does not provide an onboard sine-wave generator. A controller
must compute each sample and stream it over CAN.

For operation mode, send Type~1 frames with
\[
  t_{ff}(t)=\tau_0\sin(2\pi f t), \qquad K_p=0,\quad K_d=0.
\]

For current mode, write
\[
  i_q(t)=\frac{\tau_0}{K_t}\sin(2\pi f t)
\]
to \texttt{iq\_ref} at index \texttt{0x7006}. For example, a
2.87\,N$\cdot$m torque amplitude corresponds to about 3.05\,A using
$K_t=0.94$\,N$\cdot$m/A.

\section{ESP32 Timing}

At 1\,Mbps CAN, a 29-bit extended frame is roughly 110 bits, or about
110\,$\mu$s on the bus. Two motors commanded at 500\,Hz each require about
1000 frames/s, so the command traffic is roughly 11\% bus utilisation before
feedback frames. A 1\,ms hardware-timer loop on the ESP32 is adequate if WiFi
and Bluetooth are disabled and the CAN/timer work is kept off latency-sensitive
tasks.

\end{document}
